% CANTERA CONSERVADA: NO INCORPORAR AUTOMÁTICAMENTE A II NI A III.
% La precisión autoral posterior limita la inserción al puente constitutivo--LC.
% Para el PDF usar únicamente insercion_focal_constitutiva_LC.tex.
% Este desarrollo extenso se conserva para futuras selecciones autorizadas.
\section{Unidad radional y transporte de fase, acción y respuesta constitutiva}
\label{iii:sec:radion}

La respuesta constitutiva permite volver desde las magnitudes del vacío hasta
los canales que las producen. Esta inversión adquiere mayor contenido al
conservar también la unidad de fase, la forma de la elipse de acción y la
memoria del recorrido. El radión reúne esas lecturas sobre el mismo estado;
su incorporación precisa qué información se transporta entre ellas y qué
normalización utiliza cada una.

El antecedente es la sección compatible $X_\infty$ construida por
APP--TRIT--TPK y residente en la estructura discreta conjunta del continuo.
APP conserva las evaluaciones aditiva y multiplicativa con sus residuos,
cocientes y hojas; TRIT selecciona el régimen y su orientación; TPK compone
selección, transporte y actualización del registro. La holonomía $\Gamma_9$
actúa sobre ese registro antes de su evaluación reducida $K_{\rm ph}$.
Las coordenadas $\pi,e,\varphi,\alpha$ y la constante reducida de Planck
$\hbar$ utilizados aquí son resultados anteriores de esa construcción.
El lector global $\Delta_4$ se explicita más abajo sobre esas coordenadas.
La operación presente recibe sus salidas y sus relaciones,
sin volver a elegir las cifras ni los estados que las originan.

\subsection{Densidad de dos hojas sobre el ambiente ternario}
\label{iii:subsec:radion-densidad}

El selector trítico de fase operativa tiene las fibras
\[
 H_+=\{1,4,7\},\qquad H_-=\{2,5,8\},\qquad H_0=\{3,6,9\}.
\]
El conjunto activo $H_{\rm act}=H_+\cup H_-$ contiene seis posiciones.
El sector de incidencia utiliza una posición activa y una pareja no ordenada
de posiciones activas. Por tanto, su dominio finito es
\[
 \mathcal D_6=H_{\rm act}\times\binom{H_{\rm act}}2,\qquad
 |\mathcal D_6|=6\binom62=90.
\]
La extensión a ocho posiciones mantiene la misma colección de parejas y
produce $N_8=8\binom62=120$. Son los dos sectores de la respuesta
\eqref{iii:eq:vacuum-op}. Su cociente $90/120$ y la densidad que se define
a continuación tienen dominios y normalizaciones diferentes.

\begin{proposition}[Densidad orientada del sector activo]
\label{iii:prop:radion-densidad}
Las dos hojas conjugadas del sector $\mathcal D_6$, normalizadas por el
ambiente ternario $\mathbb F_3^6$, determinan la densidad
\begin{equation}
 \rho_R=\frac{|\{+,-\}\times\mathcal D_6|}{|\mathbb F_3^6|}
       =\frac{2\cdot90}{729}=\frac{20}{81}.
 \label{iii:eq:radion-rho}
\end{equation}
\end{proposition}
\begin{proof}
La proyección sobre la etiqueta de hoja tiene dos fibras disjuntas, cada
una de cardinal $90$. La suma de sus cardinales es $180$; la normalización
por $3^6=729$ y la reducción por $9$ dan $20/81$.
\end{proof}

El denominador es el ambiente de $729$ palabras ternarias, no el conjunto
de $243$ palabras visibles ni las $468$ emisiones decimales del censo
inicial. La densidad conserva ambos ejemplares orientados del sector.
La restricción a una sola hoja produciría $10/81$ y, por ello, una sección
torsional distinta.

\subsection{Secciones directa y torsional de la unidad de fase}
\label{iii:subsec:radion-secciones}

La vuelta construida tiene coordenada $T_+=2\pi$. En la carta angular del
registro dodecafásico, los centros y el origen de fase son
\begin{equation}
 \theta_m^\circ=(30m-10)^\circ,\qquad 1\leq m\leq12.
 \label{iii:eq:radion-calendario}
\end{equation}
La segunda posición es $\theta_2^\circ=50^\circ$ y representa la fracción
$5/36$ de la vuelta. Este enunciado conserva el origen de fase de la carta:
rotar la numeración cambia los representantes angulares. La ecuación
$30m-10=50$ tiene la única solución $m=2$ dentro del registro.

La coordenada directa $A_{\deg}=1000\alpha$ y su carta en radianes
$x=(\pi/180)A_{\deg}$ ya se han construido en
\eqref{hmtII:eq:ii-par-angular}. Al sumar la segunda fase se obtiene
\begin{equation}
 S_E=\frac{T_+}{360}(50+1000\alpha)
     =2\pi\left(\frac5{36}+\frac{25}{9}\alpha\right)
     =x+\frac{5\pi}{18}.
 \label{iii:eq:radion-SE}
\end{equation}
El subíndice $E$ identifica la sección directa o común. Su vínculo con las
lecturas eléctricas se establece mediante los canales y el operador
constitutivo; no se identifica el escalar $S_E$ con un campo eléctrico.

En el perfil considerado, los cilindros de las coordenadas construidas
permiten utilizar los intervalos racionales
\[
 3.14159<\pi<3.14160,\qquad
 0.00729735<\alpha<0.00729736.
\]
Como los factores de \eqref{iii:eq:radion-SE} son positivos y crecientes,
la sustitución de los extremos inferiores y superiores prueba
$1<S_E<2$. La división por la unidad determina entonces
\begin{equation}
 q_E=\lfloor S_E\rfloor=1,\qquad D_E=S_E-1.
 \label{iii:eq:radion-DE}
\end{equation}
Este cociente clasifica una celda de la recta de fase. El número de retornos
conservado en la memoria pertenece a otro lector.

La segunda sección utiliza el defecto global construido sobre los lectores
regionales de $\pi$ y $e$:
\begin{equation}
 \Delta_4=\frac{\pi-e\log\pi}{270}
             \left(1+\frac{\pi}{729}\right),\qquad
 \Phi_T=\rho_R\Delta_4,\qquad S_T=1+\Phi_T.
 \label{iii:eq:radion-ST}
\end{equation}
Aquí $729=3^6$ conserva la normalización ambiente; $270$ es el
complemento $999-729$ en la cámara de base $1000$, con la unidad de cierre
separada. La operación que define $\Delta_4$ se conserva como el lector
global anterior; contar el complemento explica el coeficiente, mientras
la fórmula completa especifica la operación en que interviene.
La sección $S_T$ se determina antes de compararla con $S_E$.

\subsection{Transición orientada y clausura exacta}
\label{iii:subsec:radion-cociclo}

\begin{definition}[Coordenada de transición del radión]
El residuo real de transición entre las secciones directa y torsional es
\begin{equation}
 \varepsilon_R^{(1)}=S_T-S_E=\Phi_T-D_E.
 \label{iii:eq:radion-epsilon}
\end{equation}
Su dominio conserva las dos secciones, la carta unitaria y su orientación.
Su coordenada en grados se obtiene mediante
$\varepsilon_R^\circ=(180/\pi)\varepsilon_R^{(1)}$.
\end{definition}

La diferencia define un cociclo aditivo de transición: para secciones
$S_a,S_b,S_c$ en una misma carta, si $c_{ab}=S_b-S_a$, entonces
$c_{ab}+c_{bc}=c_{ac}$ y $c_{ba}=-c_{ab}$. Ambas identidades se siguen
por cancelación de la sección intermedia. El cociclo real entre secciones
y el cociente entero de una división posicional son objetos distintos;
este último permite evaluar el primero por bloques.

\begin{theorem}[Clausura unitaria y transporte a la carta angular]
\label{iii:thm:radion-cierre}
Las secciones \eqref{iii:eq:radion-SE} y \eqref{iii:eq:radion-ST}, con la
transición \eqref{iii:eq:radion-epsilon}, satisfacen
\begin{equation}
 1=S_E-\rho_R\Delta_4+\varepsilon_R^{(1)}.
 \label{iii:eq:radion-unidad}
\end{equation}
Si $\Delta_4^\circ=(180/\pi)\Delta_4$, la misma igualdad se expresa como
\begin{equation}
 \frac{180^\circ}{\pi}
 =50^\circ+A_{\deg}\,{}^\circ
   -\frac{20}{81}\Delta_4^\circ+\varepsilon_R^\circ.
 \label{iii:eq:radion-grados}
\end{equation}
\end{theorem}
\begin{proof}
Sustituir $\varepsilon_R^{(1)}=\rho_R\Delta_4-(S_E-1)$ en el segundo
miembro de \eqref{iii:eq:radion-unidad} da $1$. Multiplicar por $180/\pi$
y usar \eqref{iii:eq:radion-SE} da \eqref{iii:eq:radion-grados}.
\end{proof}

El contenido constructivo reside en los operandos: fase dodecafásica,
desplazamiento de base $1000$, densidad de dos hojas y defecto global
se fijan antes de tomar su diferencia. La clausura es una identidad exacta
de esa diferencia. Su interpretación no requiere asignarle una función
correctora sobre ninguna de las dos vías de $\alpha$.

\subsection{Evaluación posicional y compatibilidad con el refinamiento}
\label{iii:subsec:radion-refinamiento}

En base $B=1000$, sean $t_i,v_i$ los bloques de dos operandos a profundidad
$N$ y sea $c_{N+1}$ el coeficiente de arrastre procedente de la continuación
retenida. La resta orientada actúa desde el extremo remoto hacia el inicio:
\begin{equation}
 u_i=t_i-v_i+c_{i+1},\qquad
 r_i=u_i-B\left\lfloor\frac{u_i}{B}\right\rfloor,\qquad
 c_i=\left\lfloor\frac{u_i}{B}\right\rfloor.
 \label{iii:eq:radion-subcarry}
\end{equation}
La división euclídea implica $0\leq r_i<B$. Fijados los operandos y la
continuación, cada par $(r_i,c_i)$ es único; la inducción descendente prueba
la unicidad de toda la resta. Además, al multiplicar por $B^{-i}$ y sumar,
los términos interiores se cancelan y queda
\begin{equation}
 c_1+\sum_{i=1}^N r_iB^{-i}
 =\sum_{i=1}^N(t_i-v_i)B^{-i}+c_{N+1}B^{-N}.
 \label{iii:eq:radion-subcarry-telescopio}
\end{equation}
Esta identidad mantiene visible el efecto de la continuación. La
representación finita con $c_{N+1}=0$ calcula exactamente la diferencia
de los operandos truncados; la representación del límite conserva el
coeficiente que transmite su cola compatible.

La evaluación de la transición completa puede escribirse como el funcional
\begin{equation}
 F(p,u,a)=1+\frac{20}{81}\frac{p-u\log p}{270}
                     \left(1+\frac{p}{729}\right)
           -2p\left(\frac5{36}+\frac{25}{9}a\right).
 \label{iii:eq:radion-funcional}
\end{equation}
Sean $I_N^\pi,I_N^e,I_N^\alpha$ los cilindros compatibles de las tres
salidas HMT, con diámetros que tienden a cero. En un compacto de $p>0$,
$F$ es continuamente diferenciable. Por ello existe una constante finita
$L$ tal que, para dos ternas $z,z'$ del compacto,
\[
 |F(z)-F(z')|\leq L\|z-z'\|_1.
\]
Se sigue que el diámetro de
$F(I_N^\pi\times I_N^e\times I_N^\alpha)$ tiende a cero. Las imágenes
son encajadas por inclusión de sus dominios; su intersección contiene
exactamente $F(\pi,e,\alpha)=\varepsilon_R^{(1)}$.
Así, la evaluación converge con las coordenadas que utiliza, incluida
$\alpha$. Las tasas particulares de los cilindros se transportan por esa
cota; un testigo numérico finito no sustituye este argumento de límite.

Dos controles exactos muestran la dependencia de la construcción. Si se
conserva una sola hoja, la transición cambia en $-(10/81)\Delta_4$.
Si se cambia $m=2$ por la fase vecina $m=1$ o $m=3$, la sección directa
cambia respectivamente en $-\pi/6$ o $+\pi/6$, y la transición cambia
con el signo contrario. Estos cambios alteran los operandos antes de
evaluarlos; permiten comprobar que fase y número de hojas intervienen
efectivamente en el resultado.

\subsection{Cuatro lecturas de una unidad orientada}
\label{iii:subsec:radion-cuatro}

\begin{definition}[Unidad radional]
La unidad radional es la sección orientada de $X_\infty$ que recorre una
unidad de fase levantada $\theta=1$, conserva la etiqueta de hoja
$\sigma\in\{+1,-1\}$, el cilindro de refinamiento y la memoria de
retorno, satisface \eqref{iii:eq:radion-unidad} y tiene acción firmada
$\sigma\hbar$ bajo el lector areal.
\end{definition}

La definición recibe la elipse y sus pruebas ya expuestas en
\ref{hmtII:sec:ellipse}--\ref{hmtII:sec:ii-elipse-lc}. Su función es
reunirlas con la transición concreta entre $S_E$ y $S_T$. Para una fase
levantada general, la representación orientada es
\[
 \gamma_\sigma(\theta)
  =(a_S\cos\theta,\sigma b_S\sin\theta),\qquad a_Sb_S=2\hbar.
\]
El cambio de signo de la segunda coordenada invierte la forma de área,
de modo que la ley ya demostrada se transporta a
$\mathcal A_\sigma(\theta)=\sigma\hbar\theta$.
Una unidad de fase tiene acción $\sigma\hbar$; una vuelta tiene acción
$\sigma h$, donde $h=2\pi\hbar$.

Los cuatro lectores conservan los siguientes objetos:
\begin{enumerate}
 \item \emph{Fase circular.} El mapa
 $\theta\mapsto e^{i\sigma\theta}$ toma valores en $S^1$ y conserva
 la clase angular módulo $2\pi$.
 \item \emph{Acción y forma.} El mapa $\theta\mapsto
 \gamma_\sigma(\theta)$ conserva los semiejes etiquetados, la razón
 $e^{-\eta_{\rm el}}=b_S/a_S$ y la acción firmada.
 \item \emph{Geometría interior.} El interior
 $Q^2/a_S^2+P^2/b_S^2<1$, con la distancia de Hilbert definida en el
 desarrollo geométrico, se transporta al disco de Beltrami--Klein por
 $(Q,P)\mapsto(Q/a_S,P/b_S)$. Este lector asigna una métrica al
 interior; sus puntos y distancias no son valores de acción.
 \item \emph{Memoria del retorno.} La división
 $n=9q_9(n)+r_9(n)$, con $0\leq r_9(n)<9$, conserva el par
 $(r_9(n),q_9(n))$. Su actualización es
 \[
 (r_9(n+9),q_9(n+9))=(r_9(n),q_9(n)+1).
 \]
 El residuo retorna y el cociente registra una vuelta adicional. Este
 lector se conserva junto al registro completo del TPK, no en su lugar.
\end{enumerate}

La última identidad se prueba sustituyendo la división de $n$ en $n+9$;
la unicidad euclídea fija ambas coordenadas. El retorno de fase deja así
una diferencia observable en la memoria. La unidad de origen de los
cuatro lectores no identifica sus codominios: la clase circular, la
acción, la distancia interior y el número de retornos cumplen funciones
complementarias dentro de la misma construcción.

\subsection{Reconstrucción constitutiva de la forma radional}
\label{iii:subsec:radion-puente}

El enlace con el vacío se obtiene componiendo mapas ya construidos. En el
dominio angular contractivo $x>|y|$, los canales son
$q_\pm=\exp[-(x\pm y)]\in(0,1)$. En la cámara $x>y>0$ se tiene
$q_+<q_-$. Definimos $s_\pm=q_\pm^{30}$ y conservamos la función
\[
 f(s)=\frac{1+s+s^2}{1+s+s^2+s^3},\qquad 0<s<1,
\]
de modo que $r_\pm=f(s_\pm)$. La respuesta etiquetada conserva
$r_\pm$ y sus proyectores; las dos coordenadas normalizadas
$\widehat Z,\widehat c$ los recuperan por \eqref{iii:eq:four-identities}.

\begin{proposition}[Transporte exacto entre respuesta constitutiva y forma LC]
\label{iii:prop:radion-puente}
Sean $(\widehat Z,\widehat c)$ coordenadas de la imagen constitutiva
normalizada, con etiquetas de hoja conservadas. Sean $s_\pm\in(0,1)$
las raíces únicas de
\begin{equation}
 r_\pm s_\pm^3+(r_\pm-1)(1+s_\pm+s_\pm^2)=0,\qquad
 r_+=\frac1{\sqrt{\widehat Z\widehat c}},\quad
 r_-=\sqrt{\frac{\widehat Z}{\widehat c}}.
 \label{iii:eq:radion-inversa}
\end{equation}
La forma de la elipse y la impedancia LC relativa se reconstruyen mediante
\begin{equation}
 \eta_{\rm el}=\frac12\log\frac{\log s_+}{\log s_-},\qquad
 \frac{Z_{\rm LC}}{Z_*}
   =e^{-\eta_{\rm el}}
   =\sqrt{\frac{\log s_-}{\log s_+}}.
 \label{iii:eq:radion-puente-LC}
\end{equation}
Conservando además $\hbar$, se reconstruyen los semiejes de acción;
conservando la carta $(\omega,Z_*)$, se reconstruyen $L_\eta$ y $C_\eta$.
\end{proposition}
\begin{proof}
La inversión cúbica de \eqref{iii:eq:cubic} da $s_\pm$ de manera
unívoca. Como $\log s_+=-30(x+y)$ y $\log s_-=-30(x-y)$, ambos logaritmos
son negativos y su cociente es positivo. Por tanto,
\[
 \frac12\log\frac{\log s_+}{\log s_-}
  =\frac12\log\frac{x+y}{x-y}
  =\operatorname{artanh}(y/x)=\eta_{\rm el}.
\]
La ley $Z_{\rm LC}/Z_*=e^{-\eta_{\rm el}}$ fue demostrada en
\eqref{hmtII:eq:ii-elipse-impedancia}. Sustituirla da la segunda
identidad. Las fórmulas de semiejes y de $L_\eta,C_\eta$ ya expuestas
completan la reconstrucción en sus respectivas cartas.
\end{proof}

El dominio de la proposición es la imagen de la respuesta normalizada:
los dos valores $r_\pm$ recuperados deben pertenecer a $(3/4,1)$.
Si las coordenadas se reciben en otra carta dimensional, primero se
deshace la transformación de \eqref{iii:eq:dimensions}. La raíz cúbica
se aplica después a las coordenadas internas. La constante de referencia
$Z_*$ y el reloj $\omega$ son datos de la realización LC; la respuesta
adimensional no reconstruye sus unidades por sí sola.

La distinción entre ambos cocientes queda expresada por sus funciones:
\begin{equation}
 \widehat Z=\frac{f(s_-)}{f(s_+)},\qquad
 \frac{Z_{\rm LC}}{Z_*}
       =\sqrt{\frac{\log s_-}{\log s_+}}.
 \label{iii:eq:radion-dos-impedancias}
\end{equation}
Comparten los canales y admiten la composición explícita anterior, pero
son evaluaciones diferentes. La inversión de las hojas intercambia
$s_+$ y $s_-$; en consecuencia invierte ambos cocientes, cambia
$\eta_{\rm el}$ por $-\eta_{\rm el}$ y conserva $\widehat c$.
En $y=0$ los canales coinciden y ambas razones valen $1$.

Finalmente, la recuperación angular
$x=-\log(q_+q_-)/2$ enlaza la respuesta con la sección directa mediante
$S_E=x+5\pi/18$. Si se conserva también la sección regional de
$\pi,e$ que determina $\Delta_4$, se recupera
\[
 \varepsilon_R^{(1)}
 =1+\frac{20}{81}\Delta_4-x-\frac{5\pi}{18}.
\]
Esta última composición indica exactamente qué transmite el vacío y qué
permanece en el antecedente regional. Las dos coordenadas constitutivas
recuperan el par angular y la forma; el estado HMT conserva además la
acción, las secciones regionales, la orientación y la memoria necesarias
para reconstruir la unidad radional completa.
